\documentclass{article}
\usepackage[T1]{fontenc}
\usepackage{lmodern}
\usepackage{graphicx}
\usepackage{amsmath,amssymb}
\usepackage[a4paper,margin=2cm]{geometry}
\usepackage[absolute,overlay]{textpos}
\setlength{\TPHorizModule}{1mm}
\setlength{\TPVertModule}{1mm}
\usepackage[indonesian]{babel}
\usepackage{hyperref}
\setlength{\emergencystretch}{2em}
% unit: Halaman 25

\begin{document}

% === Halaman 25 (python/native) ===

3. Distinguishing Attack

• Tujuan: menentukan apakah suatu data merupakan keluaran algoritma kriptografi 

atau data yang benar-benar acak (random).

• Jika secara statistik penyerang dapat membedakannya dengan probabilitas yang 

secara signifikan lebih baik daripada tebakan acak, berarti terdapat distinguishing 

attack.

• Contoh: misalkan sebuah cipher menghasilkan:

               101010101010101010101010...

    Pola tersebut jelas tidak menyerupai random. 

    Penyerang dapat membuat distinguisher sederhana: 

    Jika jumlah 0 dan 1 sangat tidak seimbang atau terdapat pola tertentu, 

kemungkinan data tersebut berasal dari cipher.

• Cipher modern dirancang agar kelemahan seperti ini tidak mudah ditemukan.

25

\end{document}
